\thispagestyle{plain}
\phantomsection
\addcontentsline{toc}{chapter}{Abstract}
\centerline{\zihao{-3}\heiti Abstract}

\linespread{1.3}\zihao{-4}\bigskip

This report addresses the 2021 CUMCM Problem A on FAST active reflector shape adjustment. We establish an integrated framework including ideal paraboloid determination, constrained node adjustment, and receiving-ratio evaluation. Under actuator stroke limits and cable-length variation constraints, the adjusted reflector is optimized to approximate the ideal working paraboloid.

The reflector region within the 300 m aperture is discretized by cable-net nodes. A weighted least-squares fitting objective is built and solved with constrained optimization methods to obtain node coordinates and actuator radial displacements. Then, geometric ray-tracing is used to compute the receiving ratio of the feed cabin effective region and compare it with the baseline spherical reflector.

The study shows that the task can be formulated as a large-scale geometric constrained optimization problem with clear physical interpretation and practical outputs. The model can directly provide the required vertex coordinates, adjusted node data, actuator displacements, and receiving-ratio metrics, and can be extended to incorporate more engineering factors.

\bigskip
\noindent{\zihao{4}\heiti Keywords:}
FAST; active reflector; paraboloid fitting; constrained optimization; receiving ratio
